% Generated draft: requires mathematical and source review.
\documentclass[11pt]{article}
\usepackage[margin=25mm]{geometry}
\usepackage{fontspec}
\setmainfont{Latin Modern Roman}
\usepackage{amsmath,amssymb,mathtools,unicode-math}
\setmathfont{Latin Modern Math}
\usepackage{enumitem,needspace,xurl,hyperref}
\setlength{\parindent}{0pt}
\setlength{\parskip}{5pt}
\setlength{\emergencystretch}{3em}
\newcommand{\N}{\mathbb{N}}
\newcommand{\Z}{\mathbb{Z}}
\newcommand{\Q}{\mathbb{Q}}
\newcommand{\R}{\mathbb{R}}
\newcommand{\C}{\mathbb{C}}
\newcommand{\F}{\mathbb{F}}
\newcommand{\A}{\mathbb{A}}
\newcommand{\PP}{\mathbb{P}}
\newcommand{\E}{\mathbb{E}}
\newcommand{\eps}{\varepsilon}
\newcommand{\dd}{\,\mathrm{d}}
\begin{document}
\section*{UnsolvedMath ID: 701765}
\textbf{Source problem code:} catalogue entry
\subsection{Definitions and mathematical statement}
For a group \(G\), use the commutator convention \([a,b]=a^{-1}b^{-1}ab\), and define iterated commutators by \([a,b,b]=[[a,b],b]\) and \([a,b,b,b]=[[[a,b],b],b]\). An element \(x\in G\) is left \(3\)-Engel if
\[
[g,x,x,x]=1\qquad\text{for every }g\in G.
\]
Write
\[
x^G=\{g^{-1}xg:g\in G\},\qquad N=\langle x^G\rangle,
\]
so that \(N\) is the normal closure of \(x\) in \(G\). A group is locally finite if every subgroup generated by finitely many elements is finite.

The supplied problem asks: for every group \(G\) of exponent \(8\) and every left \(3\)-Engel element \(x\in G\) of order exactly \(2\), must \(N=\langle x^G\rangle\) be locally finite? Equivalently, under these hypotheses, must
\[
\bigl|\langle h_1,\ldots,h_r\rangle\bigr|<\infty
\qquad\text{for every integer }r\geq1\text{ and every }h_1,\ldots,h_r\in N?
\]
Here \(x^2=1\) and \(x\ne1\). The exponent hypothesis implies \(g^8=1\) for every \(g\in G\). The source wording is retained: it does not explicitly distinguish exact exponent \(8\) from the convention that ``exponent \(8\)'' means that the exponent divides \(8\). This convention requires confirmation in mathematical review.
\subsection{Short English statement}
G. Traustason asks whether the normal closure of every involution that is left \(3\)-Engel in a group of exponent \(8\) is locally finite. The supplied record labels the problem open; this draft does not independently verify that status.
\subsection{Sources}
\begin{itemize}
\item Local UnsolvedMath registry, supplied complete source record: permanent ID \(701765\), exact source code \texttt{catalogue entry}, title \emph{Kourovka Notebook Problem 20.107}; upstream revision \texttt{d371ee917ac9647fbb06cd6fc529deec6f3dda82}. The supplied status is \texttt{open}.
\item Supplied source attribution: \emph{The Kourovka Notebook (Unsolved Problems in Group Theory)}, 20th issue (2022), problem 20.107; also identified with the 21st issue edition (2026), arXiv:1401.0300v48. Proposed by G. Traustason in 2022. The supplied collection description names E. I. Khukhro and V. D. Mazurov as editors. These bibliographic details are supplied attribution and have not been independently checked.
\end{itemize}
\end{document}
